\begin{document}
\MakeTitlePage
\pagenumbering{gobble}
\setcounter{page}{2}
\includepdf[pages=-]{../assignments/prochazkaelias23-zadani.pdf}
\includepdf[pages=-]{../assignments/nagyfilip23-zadani.pdf}

\section*{Anotace}
Práce se zabývá návrhem modulárního školního informačního systému bod se
zaměřením na správu rozvrhu, učeben a výukových skupin. Zahrnuje návrh
uživatelského rozhraní, backendu a databáze. Dokument je pracovní osnovou; popis implementace,
výsledky testování a závěry budou doplňovány podle skutečného stavu projektu.

\section*{Klíčová slova}
školní informační systém, rozvrh, výukové skupiny, učebny, webová aplikace

\section*{Annotation}
The project concerns the design of bod, a modular school information system
focused on timetables, classrooms and teaching groups. It covers the design
of the user interface, backend and database. This document is a working outline;
implementation details, test results and conclusions will be added as the
project progresses.

\section*{Key words}
school information system, timetable, teaching groups, classrooms, web application

\newpage
\section*{Prohlášení}
\begin{citacePlain}
Prohlašujeme, že jsme tuto práci vypracovali společně a že byly využity pouze
literární prameny a informační zdroje uvedené v seznamu použité literatury a zdrojů.

Prohlašujeme, že při zpracování práce byly nástroje umělé inteligence využity
v souladu s principy akademické integrity a že jejich použití je v práci
vhodným způsobem uvedeno.

Prohlašujeme, že odevzdané verze, jedna v tištěné podobě a elektronická verze
na GitLabu jsou stejné.

Souhlasíme s využitím této práce pro výukové účely VOŠ a SPŠE Plzeň.
\end{citacePlain}
\vspace{1cm}
\noindent V Plzni dne: \dotfill
\vspace{1cm}

\noindent\begin{tabular}{@{}p{0.46\textwidth}p{0.46\textwidth}@{}}
\dotfill & \dotfill \\
Eliáš Jan Procházka & Filip Nagy
\end{tabular}

\clearpage
\tableofcontents
\clearpage
\edef\FirstTextPage{\arabic{page}}
\pagenumbering{arabic}
\setcounter{page}{\FirstTextPage}

\tocchapter{Úvod}
Školní informační systém propojuje správu uživatelů, výukových skupin,
učeben a rozvrhu. Projekt bod vzniká jako týmová dlouhodobá maturitní práce.
Cílem je navrhnout a implementovat propojené uživatelské rozhraní, API
systému a datový model podle přiložených zadání.

\chapter{Analýza požadavků}
\section{Práce rozvrháře}
Kapitola bude obsahovat popis přípravy úvazků, dělení skupin a kontroly
kolizí. Poznatky z konzultací budou odlišeny od implementovaných funkcí.

\section{Rozdělení práce}
Eliáš Jan Procházka zpracovává frontend a uživatelské rozhraní
(kapitola~\ref{chap:frontend}). Filip Nagy zpracovává backend a databázi
(kapitola~\ref{chap:backend}). Analýza požadavků a návrh architektury jsou společné.

\chapter{Návrh řešení}
Kapitola bude popisovat hranice modulů, vazbu frontendu a API a způsob
správy dat. Diagramy musí odpovídat skutečné implementaci.

\chapter{Frontend a uživatelské rozhraní}
\label{chap:frontend}
\textbf{Autor části:} Eliáš Jan Procházka.

\section{Architektura frontendu}
Zde bude popsáno uspořádání aplikace, společné komponenty a hranice modulů.

\section{Práce s rozvrhem}
Popis přípravy podkladů, sestavování a kontroly rozvrhu bude doplněn
podle skutečně implementovaných funkcí.

\section{Propojení s API}
Zde bude popsána komunikace s backendem, přihlášení a zpracování chyb.

\chapter{Backend a databáze}
\label{chap:backend}
\textbf{Autor části:} Filip Nagy.

\section{Datový model}
Zde bude popsán datový model a vazby mezi uživateli, skupinami, učebnami
a rozvrhem.

\section{Aplikační rozhraní}
Zde bude popsáno API, autentizace a vyhodnocení oprávnění.

\section{Validace rozvrhu}
Kontroly kolizí a podmínky publikace budou popsány podle skutečné implementace.


\chapter{Testování a vyhodnocení}
Kapitola bude obsahovat reprodukovatelné postupy a skutečné výsledky.
Naměřené hodnoty budou uvedeny s verzí aplikace a podmínkami měření.

\tocchapter{Závěr}
Závěr bude doplněn po vyhodnocení stanovených cílů. Pracovní osnova
nepředstavuje tvrzení o dokončení zadání.

\printbibliography[heading=bibintoc]
\tocchapter{Seznam zkratek}
\begin{description}
  \item[API] Aplikační programové rozhraní.
  \item[CI/CD] Průběžná integrace a nasazování.
  \item[IS] Informační systém.
\end{description}

\tocchapter{Přílohy}
\StartAppendices
\chapter{Použité nástroje}
Dokument je sestaven ze školní šablony pro dlouhodobé maturitní práce.
K inicializaci struktury dokumentace a sestavovacích skriptů byl použit
asistent OpenAI Codex. Odborný text, zdroje a výsledky musí ověřit autoři.
Další využití asistivních nástrojů bude popsáno podle skutečné práce autorů.
\end{document}
